\chapter{Conclusion}

The main goals of this thesis were to provide an introduction to spectral graph theory and to give deeper insight into two recent research topics.
First, we presented the very foundations of graph theory, collected useful results from linear algebra and applied them to prove elementary facts about graph spectra.
Chapter 3 was devoted to classical results, establishing the connection between the structural and spectral properties of a graph.
We introduced the concept of graph expansion and showed that a high expansion ratio is qualitatively equivalent to a large spectral gap (Cheeger's inequality).

The first major application we discussed was the recent work of Jiang, Tidor, Yao, Zhang, and Zhao \cite{equiangular}, in which the authors determined the number of equiangular lines with a fixed angle in high dimensions. 
To reduce this geometric question to a graph-theoretic problem, they considered a graph associated with the Gram matrix of a set of unit vectors representing the equiangular set.
Then, they used the knowledge of forbidden induced subgraphs together with a switching argument to show that this graph can be assumed to have bounded degree.
Their main contribution to spectral graph theory was a novel result on the distribution of the eigenvalues of bounded-degree graphs: 
The second eigenvalue multiplicity of a connected $n$-vertex graph with bounded maximum degree is at most $O(n / \log \log n)$.
Although \cite{equiangular} only needed the fact that the multiplicity is sublinear, they raised the question of whether their bound could be improved.

Addressing this question, we established stronger upper bounds on the second eigenvalue multiplicity for graphs with small balanced separators in Theorem \ref{thm:separable-snd-eigval-mult}.
As a corollary, we obtained that graphs with a fixed excluded minor have second eigenvalue multiplicity $O(\sqrt{n})$.
It would be interesting to know whether there are graph classes with small separators for which the bounds of Theorem \ref{thm:separable-snd-eigval-mult} are tight. 
Besides that, it might be helpful for further research to find other sufficient conditions for graphs to have small second eigenvalue multiplicity.

The second recent research result we discussed was the proof of the Sensitivity Conjecture. 
Building on a reduction due to Gotsman and Linial, Huang \cite{sensitivity} resolved this long-standing conjecture by proving that every induced subgraph on more than half of the vertices of the $n$-dimensional hypercube has maximum degree at least $\sqrt{n}$.
Huang constructed a particular signed adjacency matrix of the hypercube and used interlacing to lower bound the largest eigenvalue of a large induced subgraph. 
This, in turn, yielded the lower bound on the maximum degree.

Spectral graph theory is a highly active field of research, of which we could merely scratch the surface. 
For further reading on the general subject, we refer the reader to \cite{brouwerhaemers}, \cite{spielman}, and \cite{godsilroyle}. 

There are numerous other interesting applications of spectral graph theory that are beyond the scope of this thesis.
One example is spectral graph drawing. 
Intuitively, a "nice" drawing of a graph in the plane is one in which neighbouring vertices are drawn close to each other, and there are few edge crossings.
An approach to graph drawing proposed by Hall \cite{hall} uses the entries of eigenvectors corresponding to the smallest nonzero Laplace eigenvalues as coordinates for the vertices.
The theory behind Hall's approach is also explained in \cite{spielman}.

Furthermore, there are profound connections between spectral graph theory and the theory of random walks on graphs.
A fundamental result in this context states that every simple random walk on a finite graph converges to a limiting stationary distribution. 
For expander graphs, that is, graphs with a large spectral gap, this happens as rapidly as possible.
An important consequence is that the set of vertices visited by a "short" random walk on an expander graph resembles (in some sense) a subset of vertices sampled independently from a uniform distribution.
The surveys \cite{randomwalksongraphs} and \cite[Section 3]{expandergraphs} provide an introduction to this topic.